% TOP_PROBLEM: {"id": 2308006, "title": "Research Problems in Function Theory — Problem 8.6", "category_id": 8, "category": {"id": 8, "name": "analysis", "display_name": "Analysis", "description": "Limits, continuity, calculus, and function theory.", "slug": "analysis", "order_index": 8, "created_at": "2026-07-31T15:26:25.671Z"}, "status": "open"}
% FIRST_POSTED: null
% ENTRY_KIND: statement_only
% Imported from: batch06.tex; UnsolvedMath ID 2308006
% Compile twice: lualatex 2308006.tex
\documentclass[11pt]{article}
\usepackage{fontspec}
\setmainfont{Latin Modern Roman}
\usepackage{amsmath,amssymb,mathtools,unicode-math}
\setmathfont{Latin Modern Math}
\usepackage{xurl,hyperref}
\hypersetup{hidelinks,pdftitle={UnsolvedMath Problem 2308006}}
\newcommand{\sref}[2]{\hyperlink{source:#1:#2}{[S#2]}}
\newcommand{\cataloguescope}[1]{\paragraph{Mathematical question.}#1}
\begin{document}
% UnsolvedMath ID 2308006; source catalogue code AMR-022-8006
\setcounter{section}{2308005}
\section{Research Problems in Function Theory — Problem 8.6}\label{2308006}
{\small\sffamily UnsolvedMath ID 2308006 · Analysis}
\subsection{Definitions and mathematical statement}

\paragraph{Shared notation.}$\mathbb N=\{1,2,\ldots\}$, and $\mathbb Z,\mathbb Q,\mathbb R,\mathbb C$ denote the integers, rationals, reals and complex numbers. Any other conventions are specified locally.


For $p>0$, the Bergman space $A^p$ consists of holomorphic $f$ on $\mathbb D$ with $\int_{\mathbb D}|f|^p\,dA<\infty$. A zero divisor is the multiset of zeros of a nonzero such function, with multiplicities and no accumulation inside the disc. Any finite multiplicity at $0$ is separated out before taking reciprocal products; write the remaining zeros $z_k\ne0$ in nondecreasing modulus. For a finite nonempty $F\subset\mathbb T$ let $I_\nu$ be the complementary open arcs and $|I_\nu|$ their angular lengths. Define
\[R(I)=\{re^{i\theta}:1-|I|/(2\pi)\le r<1,\ e^{i\theta}\in\overline I\},
\quad G_F=\mathbb D\setminus\bigcup_\nu R(I_\nu),
\]
\[\kappa(F)=\frac1{2\pi}\sum_\nu\frac{|I_\nu|}{2\pi}
\log\big(2\pi/|I_\nu|+1\big).
\]
Sums over zeros in $G_F$ count multiplicity.
\paragraph{Statement recorded in the supplied source.}
Characterize zero divisors for each fixed $A^p$, or at least give a nontrivial converse to the necessary estimate
\[\prod_{k=1}^N|z_k|^{-1}=O(N^{1/p+\varepsilon})\quad(N\to\infty)
\quad\text{for every }\varepsilon>0,
\]
whose exponent cannot generally omit $\varepsilon$. The source update supplies a modified Korenblum criterion: for an absolute $\lambda>1$, finiteness of
\[\sup_F\big\{\sum_{z_k\in G_F}(1-|z_k|)-\lambda p^{-1}\kappa(F)\big\}
\]
is necessary, while finiteness with $\lambda$ replaced by $\lambda^{-1}$ is sufficient. It also states that
\[\sup_F\kappa(F)^{-1}\sum_{z_k\in G_F}(1-|z_k|)<\infty
\]
characterizes being a zero divisor for some $A^p$, $p>0$. These two nearby fixed-$p$ bounds do not replace the request for an exact fixed-$p$ characterization. \sref{2308006}{2}
\subsection{Short English statement}
Characterize Bergman-space zero sequences at each exponent, or find meaningful converse conditions to the reciprocal-product bound, preserving the source’s necessary and sufficient comparison criteria.
\subsection{Sources}
\begin{description}
\item[\hypertarget{source:2308006:1}{S1}] ULAM.AI and the UnsolvedMath Contributors, UnsolvedMath: A Curated Collection of Open Mathematics Problems, record 2308006 (AMR-022-8006). CC BY 4.0. \url{https://huggingface.co/datasets/ulamai/UnsolvedMath}.
\item[\hypertarget{source:2308006:2}{S2}] Walter K. Hayman and Edward F. Lingham, Research Problems in Function Theory (2018), Problem 8.6, its complete update and relevant preceding definitions. \url{https://arxiv.org/abs/1809.07200}.
\end{description}
\label{end:2308006}
\end{document}
